\section{Completeness: classical results and the open problem of this thesis}
There is an important distinction to make at the outset. The completeness of the \emph{classical}
Lambek calculus is not an open problem in general. Several major completeness results are
already known. \citet{Dosen1985} proved a completeness theorem for the Lambek calculus of
syntactic categories; \citet{AndrekaMikulas1994} established strong completeness results for
relational semantics; and \citet{Pentus1995} proved completeness with respect to language
models. In addition, Pentus' work on the context-free status of Lambek grammars clarified the
place of the calculus in formal-language theory \citep{Pentus1993,Pentus1997}.

The open problem relevant to the present thesis is therefore narrower and more specific:
\emph{is the probabilistic semantics introduced here complete for the Lambek calculus?} This
question is open because the semantics developed in this thesis is corpus-relative, numerical,
and intentionally partial in the sense that it is tied to observed linguistic data. The standard
canonical-model and algebraic techniques from classical completeness proofs do not transfer
immediately to such a setting.

For that reason the present thesis focuses on soundness first. Establishing soundness is the right
first milestone: before asking whether the semantics captures \emph{all} valid inferences, one
must first show that every derivable Lambek inequality is semantically valid. Completeness is the
next goal, not a claim of the present work.

